\section{Exercices d'application}

\rubrique{Formules d'addition et de duplication}

\exo{}\diff{1}
Sachant que $\cos\frac{\pi}{3} = \frac12$ et $\sin\frac{\pi}{3} = \frac{\sqrt3}{2}$, calculer les valeurs exactes de $\cos\frac{\pi}{12}$ et $\sin\frac{\pi}{12}$ en utilisant les formules d'addition.

\exo{}\diff{1}
On donne $\cos a = \frac35$ avec $0 < a < \frac{\pi}{2}$ et $\sin b = \frac{5}{13}$ avec $\frac{\pi}{2} < b < \pi$.
\begin{enumerate}[label=\alph*)]
    \item Calculer $\sin a$ et $\cos b$.
    \item En déduire $\cos(a+b)$ et $\sin(a-b)$.
\end{enumerate}

\exo{}\diff{2}
Démontrer les identités suivantes :
\begin{enumerate}[label=\alph*)]
    \item $\cos\left(x+\frac{\pi}{3}\right) + \cos\left(x-\frac{\pi}{3}\right) = \cos x$
    \item $\sin\left(\frac{\pi}{4}+x\right) - \sin\left(\frac{\pi}{4}-x\right) = \sqrt2\sin x$
\end{enumerate}

\exo{}\diff{2}
Soit $x$ un réel tel que $\sin x = \frac{4}{5}$ et $\frac{\pi}{2} < x < \pi$.
\begin{enumerate}[label=\alph*)]
    \item Calculer $\cos x$.
    \item En utilisant les formules de duplication, calculer $\sin(2x)$ et $\cos(2x)$.
\end{enumerate}

\rubrique{Transformation de $a\cos x + b\sin x$}

\exo{}\diff{2}
Mettre sous la forme $R\cos(x-\varphi)$ (avec $R>0$ et $\varphi\in[0,2\pi[$) les expressions suivantes :
\begin{enumerate}[label=\alph*)]
    \item $f(x) = \cos x + \sqrt3\sin x$
    \item $g(x) = 3\cos x - 4\sin x$
\end{enumerate}

\exo{}\diff{2}
Résoudre dans $\R$ l'équation $\cos x + \sin x = 1$ en utilisant la transformation $a\cos x + b\sin x$.

\rubrique{Équations trigonométriques}

\exo{}\diff{1}
Résoudre dans $\R$ les équations suivantes :
\begin{enumerate}[label=\alph*)]
    \item $\cos x = \frac{\sqrt2}{2}$
    \item $\sin x = -\frac12$
    \item $\tan x = \sqrt3$
\end{enumerate}

\exo{}\diff{2}
Résoudre dans $[0,2\pi[$ les équations :
\begin{enumerate}[label=\alph*)]
    \item $\cos(2x) = \frac12$
    \item $\sin\left(x+\frac{\pi}{4}\right) = \frac{\sqrt3}{2}$
\end{enumerate}

\exo{}\diff{2}
Résoudre dans $\R$ l'équation $2\cos^2 x - 3\cos x + 1 = 0$.

\exo{}\diff{3}
Résoudre dans $[-\pi,\pi]$ l'équation $\sin(2x) = \cos x$.

\rubrique{Inéquations trigonométriques}

\exo{}\diff{2}
Résoudre dans $[0,2\pi[$ l'inéquation $\cos x \ge \frac12$.

\exo{}\diff{3}
Résoudre dans $]-\pi,\pi]$ l'inéquation $\sin x > \frac{\sqrt2}{2}$.

\section{Exercices d'entraînement}

\rubrique{Formules d'addition et de duplication}

\exo{}\diff{1}
Calculer $\cos\frac{7\pi}{12}$ et $\sin\frac{7\pi}{12}$ en utilisant $\frac{\pi}{3}$ et $\frac{\pi}{4}$.

\exo{}\diff{1}
Simplifier $A = \cos(a+b)\cos(a-b) - \sin(a+b)\sin(a-b)$.

\exo{}\diff{2}
Démontrer que $\cos^4 x - \sin^4 x = \cos(2x)$.

\exo{}\diff{2}
Soit $x$ un réel tel que $\cos x = -\frac{3}{5}$ et $\pi < x < \frac{3\pi}{2}$. Calculer $\sin x$, $\sin(2x)$ et $\cos(2x)$.

\exo{}\diff{2}
Exprimer $\sin(3x)$ en fonction de $\sin x$ uniquement.

\exo{}\diff{2}
Démontrer que $\frac{\sin(2x)}{1+\cos(2x)} = \tan x$ pour $x \neq \frac{\pi}{2} + k\pi$.

\rubrique{Transformation de $a\cos x + b\sin x$}

\exo{}\diff{2}
Mettre sous la forme $R\sin(x+\varphi)$ les expressions :
\begin{enumerate}[label=\alph*)]
    \item $h(x) = \sqrt3\cos x + \sin x$
    \item $k(x) = \cos x - \sin x$
\end{enumerate}

\exo{}\diff{2}
Résoudre dans $\R$ l'équation $\sqrt3\cos x - \sin x = \sqrt2$.

\exo{}\diff{3}
Soit $f(x) = 2\cos x + 3\sin x$. Déterminer le maximum et le minimum de $f$ sur $\R$.

\rubrique{Équations trigonométriques}

\exo{}\diff{1}
Résoudre dans $\R$ :
\begin{enumerate}[label=\alph*)]
    \item $\cos x = -1$
    \item $\sin x = 0$
    \item $\tan x = -1$
\end{enumerate}

\exo{}\diff{2}
Résoudre dans $[0,2\pi[$ :
\begin{enumerate}[label=\alph*)]
    \item $\sin(3x) = \frac{\sqrt3}{2}$
    \item $\cos\left(2x-\frac{\pi}{6}\right) = -\frac12$
\end{enumerate}

\exo{}\diff{2}
Résoudre dans $\R$ l'équation $2\sin^2 x - \sin x - 1 = 0$.

\exo{}\diff{2}
Résoudre dans $[0,\pi]$ l'équation $\cos(2x) + \cos x = 0$.

\exo{}\diff{3}
Résoudre dans $\R$ l'équation $\sin x + \sin(2x) = 0$.

\rubrique{Inéquations trigonométriques}

\exo{}\diff{2}
Résoudre dans $[0,2\pi[$ l'inéquation $\sin x \le \frac12$.

\exo{}\diff{2}
Résoudre dans $]-\pi,\pi]$ l'inéquation $\tan x \ge 1$.

\exo{}\diff{3}
Résoudre dans $[0,2\pi[$ l'inéquation $2\cos^2 x - \cos x - 1 < 0$.

\section{Problèmes type Probatoire/Bac}

\sujetbac[Probatoire blanc — Yaoundé]
\noindent\textbf{Partie A}
\begin{enumerate}[label=\arabic*)]
    \item Résoudre dans $\R$ l'équation $2\cos^2 x - 3\cos x + 1 = 0$.
    \item En déduire les solutions dans $[0,2\pi[$.
\end{enumerate}

\textbf{Partie B}
\begin{enumerate}[label=\arabic*)]
    \item Mettre $f(x) = \cos x + \sqrt3\sin x$ sous la forme $R\cos(x-\varphi)$.
    \item Résoudre dans $\R$ l'équation $f(x) = \sqrt2$.
    \item Résoudre dans $[0,2\pi[$ l'inéquation $f(x) \ge 0$.
\end{enumerate}

\sujetbac[Baccalauréat blanc — Douala]
\noindent\textbf{Partie A}
Soit $x$ un réel tel que $\sin x = \frac{3}{5}$ et $\frac{\pi}{2} < x < \pi$.
\begin{enumerate}[label=\arabic*)]
    \item Calculer $\cos x$.
    \item Calculer $\sin(2x)$ et $\cos(2x)$.
    \item En déduire $\tan(2x)$.
\end{enumerate}

\textbf{Partie B}
\begin{enumerate}[label=\arabic*)]
    \item Résoudre dans $\R$ l'équation $\sin(2x) = \frac{24}{25}$.
    \item Donner les solutions dans $[0,\pi]$.
\end{enumerate}

\section{Corrigés}

\corrige{de l'exercice 1}
On utilise $\frac{\pi}{12} = \frac{\pi}{3} - \frac{\pi}{4}$.
\[
\cos\frac{\pi}{12} = \cos\left(\frac{\pi}{3} - \frac{\pi}{4}\right) = \cos\frac{\pi}{3}\cos\frac{\pi}{4} + \sin\frac{\pi}{3}\sin\frac{\pi}{4} = \frac12\times\frac{\sqrt2}{2} + \frac{\sqrt3}{2}\times\frac{\sqrt2}{2} = \frac{\sqrt2 + \sqrt6}{4}
\]
\[
\sin\frac{\pi}{12} = \sin\left(\frac{\pi}{3} - \frac{\pi}{4}\right) = \sin\frac{\pi}{3}\cos\frac{\pi}{4} - \cos\frac{\pi}{3}\sin\frac{\pi}{4} = \frac{\sqrt3}{2}\times\frac{\sqrt2}{2} - \frac12\times\frac{\sqrt2}{2} = \frac{\sqrt6 - \sqrt2}{4}
\]

\corrige{de l'exercice 2}
\begin{enumerate}[label=\alph*)]
    \item $\cos a = \frac35$ et $0<a<\frac{\pi}{2}$ donc $\sin a = \sqrt{1-\cos^2 a} = \sqrt{1-\frac{9}{25}} = \sqrt{\frac{16}{25}} = \frac45$.
    $\sin b = \frac{5}{13}$ et $\frac{\pi}{2}<b<\pi$ donc $\cos b = -\sqrt{1-\sin^2 b} = -\sqrt{1-\frac{25}{169}} = -\sqrt{\frac{144}{169}} = -\frac{12}{13}$.
    \item $\cos(a+b) = \cos a\cos b - \sin a\sin b = \frac35\times\left(-\frac{12}{13}\right) - \frac45\times\frac{5}{13} = -\frac{36}{65} - \frac{20}{65} = -\frac{56}{65}$.
    $\sin(a-b) = \sin a\cos b - \cos a\sin b = \frac45\times\left(-\frac{12}{13}\right) - \frac35\times\frac{5}{13} = -\frac{48}{65} - \frac{15}{65} = -\frac{63}{65}$.
\end{enumerate}

\corrige{de l'exercice 3}
\begin{enumerate}[label=\alph*)]
    \item $\cos\left(x+\frac{\pi}{3}\right) + \cos\left(x-\frac{\pi}{3}\right) = \cos x\cos\frac{\pi}{3} - \sin x\sin\frac{\pi}{3} + \cos x\cos\frac{\pi}{3} + \sin x\sin\frac{\pi}{3} = 2\cos x\cos\frac{\pi}{3} = 2\cos x\times\frac12 = \cos x$.
    \item $\sin\left(\frac{\pi}{4}+x\right) - \sin\left(\frac{\pi}{4}-x\right) = \sin\frac{\pi}{4}\cos x + \cos\frac{\pi}{4}\sin x - (\sin\frac{\pi}{4}\cos x - \cos\frac{\pi}{4}\sin x) = 2\cos\frac{\pi}{4}\sin x = 2\times\frac{\sqrt2}{2}\sin x = \sqrt2\sin x$.
\end{enumerate}

\corrige{de l'exercice 4}
\begin{enumerate}[label=\alph*)]
    \item $\sin x = \frac45$ et $\frac{\pi}{2}<x<\pi$ donc $\cos x = -\sqrt{1-\frac{16}{25}} = -\sqrt{\frac{9}{25}} = -\frac35$.
    \item $\sin(2x) = 2\sin x\cos x = 2\times\frac45\times\left(-\frac35\right) = -\frac{24}{25}$.
    $\cos(2x) = \cos^2 x - \sin^2 x = \frac{9}{25} - \frac{16}{25} = -\frac{7}{25}$.
\end{enumerate}

\corrige{de l'exercice 5}
\begin{enumerate}[label=\alph*)]
    \item $f(x) = \cos x + \sqrt3\sin x$. $R = \sqrt{1^2 + (\sqrt3)^2} = \sqrt{1+3} = 2$.
    $\cos\varphi = \frac{1}{2}$ et $\sin\varphi = \frac{\sqrt3}{2}$ donc $\varphi = \frac{\pi}{3}$.
    Ainsi $f(x) = 2\cos\left(x-\frac{\pi}{3}\right)$.
    \item $g(x) = 3\cos x - 4\sin x$. $R = \sqrt{3^2 + (-4)^2} = \sqrt{9+16} = 5$.
    $\cos\varphi = \frac{3}{5}$ et $\sin\varphi = -\frac{4}{5}$ donc $\varphi = -\arctan\frac{4}{3}$ (mod $2\pi$).
    Ainsi $g(x) = 5\cos\left(x-\varphi\right)$ avec $\varphi = -\arctan\frac{4}{3}$.
\end{enumerate}

\corrige{de l'exercice 6}
$\cos x + \sin x = \sqrt2\cos\left(x-\frac{\pi}{4}\right)$ (car $R=\sqrt2$, $\varphi=\frac{\pi}{4}$).
L'équation devient $\sqrt2\cos\left(x-\frac{\pi}{4}\right) = 1$, soit $\cos\left(x-\frac{\pi}{4}\right) = \frac{\sqrt2}{2}$.
Donc $x-\frac{\pi}{4} = \frac{\pi}{4} + 2k\pi$ ou $x-\frac{\pi}{4} = -\frac{\pi}{4} + 2k\pi$, $k\in\Z$.
Soit $x = \frac{\pi}{2} + 2k\pi$ ou $x = 2k\pi$, $k\in\Z$.

\corrige{de l'exercice 7}
\begin{enumerate}[label=\alph*)]
    \item $\cos x = \frac{\sqrt2}{2} \iff x = \frac{\pi}{4} + 2k\pi$ ou $x = -\frac{\pi}{4} + 2k\pi$, $k\in\Z$.
    \item $\sin x = -\frac12 \iff x = -\frac{\pi}{6} + 2k\pi$ ou $x = \frac{7\pi}{6} + 2k\pi$, $k\in\Z$.
    \item $\tan x = \sqrt3 \iff x = \frac{\pi}{3} + k\pi$, $k\in\Z$.
\end{enumerate}

\corrige{de l'exercice 8}
\begin{enumerate}[label=\alph*)]
    \item $\cos(2x) = \frac12 \iff 2x = \frac{\pi}{3} + 2k\pi$ ou $2x = -\frac{\pi}{3} + 2k\pi$, $k\in\Z$.
    Soit $x = \frac{\pi}{6} + k\pi$ ou $x = -\frac{\pi}{6} + k\pi$.
    Dans $[0,2\pi[$ : $x = \frac{\pi}{6}, \frac{5\pi}{6}, \frac{7\pi}{6}, \frac{11\pi}{6}$.
    \item $\sin\left(x+\frac{\pi}{4}\right) = \frac{\sqrt3}{2} \iff x+\frac{\pi}{4} = \frac{\pi}{3} + 2k\pi$ ou $x+\frac{\pi}{4} = \frac{2\pi}{3} + 2k\pi$, $k\in\Z$.
    Soit $x = \frac{\pi}{12} + 2k\pi$ ou $x = \frac{5\pi}{12} + 2k\pi$.
    Dans $[0,2\pi[$ : $x = \frac{\pi}{12}, \frac{5\pi}{12}$.
\end{enumerate}

\corrige{de l'exercice 9}
$2\cos^2 x - 3\cos x + 1 = 0$. Posons $X = \cos x$. L'équation devient $2X^2 - 3X + 1 = 0$.
$\Delta = 9 - 8 = 1$. $X_1 = \frac{3-1}{4} = \frac12$, $X_2 = \frac{3+1}{4} = 1$.
Donc $\cos x = \frac12$ ou $\cos x = 1$.
$\cos x = \frac12 \iff x = \pm\frac{\pi}{3} + 2k\pi$, $k\in\Z$.
$\cos x = 1 \iff x = 2k\pi$, $k\in\Z$.

\corrige{de l'exercice 10}
$\sin(2x) = \cos x \iff 2\sin x\cos x - \cos x = 0 \iff \cos x(2\sin x - 1) = 0$.
Donc $\cos x = 0$ ou $\sin x = \frac12$.
$\cos x = 0 \iff x = \frac{\pi}{2} + k\pi$, $k\in\Z$.
$\sin x = \frac12 \iff x = \frac{\pi}{6} + 2k\pi$ ou $x = \frac{5\pi}{6} + 2k\pi$, $k\in\Z$.
Dans $[-\pi,\pi]$ : $x = -\frac{\pi}{2}, \frac{\pi}{2}, \frac{\pi}{6}, \frac{5\pi}{6}$.

\corrige{de l'exercice 11}
$\cos x \ge \frac12$ dans $[0,2\pi[$.
$\cos x = \frac12$ pour $x = \frac{\pi}{3}$ et $x = \frac{5\pi}{3}$.
Sur le cercle trigonométrique, $\cos x \ge \frac12$ correspond à $x \in [0,\frac{\pi}{3}] \cup [\frac{5\pi}{3},2\pi[$.

\corrige{de l'exercice 12}
$\sin x > \frac{\sqrt2}{2}$ dans $]-\pi,\pi]$.
$\sin x = \frac{\sqrt2}{2}$ pour $x = \frac{\pi}{4}$ et $x = \frac{3\pi}{4}$.
Sur le cercle, $\sin x > \frac{\sqrt2}{2}$ correspond à $x \in ]\frac{\pi}{4}, \frac{3\pi}{4}[$.

\corrige{du problème 1}
\textbf{Partie A}
\begin{enumerate}[label=\arabic*)]
    \item $2\cos^2 x - 3\cos x + 1 = 0$. Posons $X = \cos x$. $2X^2 - 3X + 1 = 0$, $\Delta = 1$, $X_1 = \frac12$, $X_2 = 1$.
    Donc $\cos x = \frac12$ ou $\cos x = 1$.
    $\cos x = \frac12 \iff x = \pm\frac{\pi}{3} + 2k\pi$, $k\in\Z$.
    $\cos x = 1 \iff x = 2k\pi$, $k\in\Z$.
    \item Dans $[0,2\pi[$ : $x = 0, \frac{\pi}{3}, \frac{5\pi}{3}$.
\end{enumerate}

\textbf{Partie B}
\begin{enumerate}[label=\arabic*)]
    \item $f(x) = \cos x + \sqrt3\sin x$. $R = \sqrt{1+3} = 2$, $\cos\varphi = \frac12$, $\sin\varphi = \frac{\sqrt3}{2}$ donc $\varphi = \frac{\pi}{3}$.
    Ainsi $f(x) = 2\cos\left(x-\frac{\pi}{3}\right)$.
    \item $f(x) = \sqrt2 \iff 2\cos\left(x-\frac{\pi}{3}\right) = \sqrt2 \iff \cos\left(x-\frac{\pi}{3}\right) = \frac{\sqrt2}{2}$.
    Donc $x-\frac{\pi}{3} = \frac{\pi}{4} + 2k\pi$ ou $x-\frac{\pi}{3} = -\frac{\pi}{4} + 2k\pi$, $k\in\Z$.
    Soit $x = \frac{7\pi}{12} + 2k\pi$ ou $x = \frac{\pi}{12} + 2k\pi$, $k\in\Z$.
    \item $f(x) \ge 0 \iff 2\cos\left(x-\frac{\pi}{3}\right) \ge 0 \iff \cos\left(x-\frac{\pi}{3}\right) \ge 0$.
    Dans $[0,2\pi[$, posons $X = x-\frac{\pi}{3}$. Alors $X \in [-\frac{\pi}{3}, \frac{5\pi}{3}[$.
    $\cos X \ge 0$ pour $X \in [-\frac{\pi}{2}, \frac{\pi}{2}] \cup [\frac{3\pi}{2}, \frac{5\pi}{3}[$.
    Donc $x \in [-\frac{\pi}{6}, \frac{5\pi}{6}] \cup [\frac{11\pi}{6}, 2\pi[$.
\end{enumerate}

\corrige{du problème 2}
\textbf{Partie A}
\begin{enumerate}[label=\arabic*)]
    \item $\sin x = \frac35$ et $\frac{\pi}{2}<x<\pi$ donc $\cos x = -\sqrt{1-\frac{9}{25}} = -\sqrt{\frac{16}{25}} = -\frac45$.
    \item $\sin(2x) = 2\sin x\cos x = 2\times\frac35\times\left(-\frac45\right) = -\frac{24}{25}$.
    $\cos(2x) = \cos^2 x - \sin^2 x = \frac{16}{25} - \frac{9}{25} = \frac{7}{25}$.
    \item $\tan(2x) = \frac{\sin(2x)}{\cos(2x)} = \frac{-\frac{24}{25}}{\frac{7}{25}} = -\frac{24}{7}$.
\end{enumerate}

\textbf{Partie B}
\begin{enumerate}[label=\arabic*)]
    \item $\sin(2x) = \frac{24}{25}$. D'après la partie A, $\sin(2x) = -\frac{24}{25}$ pour $x$ donné, mais ici on cherche les solutions générales.
    $\sin(2x) = \frac{24}{25} \iff 2x = \arcsin\frac{24}{25} + 2k\pi$ ou $2x = \pi - \arcsin\frac{24}{25} + 2k\pi$, $k\in\Z$.
    Soit $x = \frac12\arcsin\frac{24}{25} + k\pi$ ou $x = \frac{\pi}{2} - \frac12\arcsin\frac{24}{25} + k\pi$, $k\in\Z$.
    \item Dans $[0,\pi]$, les solutions sont $x = \frac12\arcsin\frac{24}{25}$ et $x = \frac{\pi}{2} - \frac12\arcsin\frac{24}{25}$.
\end{enumerate}